% Part: many-valued-logic
% Chapter: three-valued-logics
% Section: introduction

\documentclass[../../../include/open-logic-section]{subfiles}

\begin{document}

\olfileid[hi]{mvl}{thr}{int}

\olsection{परिचय}

$\True$ और $\False$ में केवल एक और मान~$\Undef$ जोड़ने पर हमें त्रिमूल्य
तर्कशास्त्र मिलता है। यद्यपि केवल एक अतिरिक्त सत्य-मान है, फिर भी $\lnot$,
$\land$, $\lor$ और $\lif$ के सत्य-फलन परिभाषित करने की संभावनाएँ काफी अधिक
हैं। तब कोई तर्कशास्त्र इन सत्य-फलनों का कोई भी संयोजन उपयोग कर सकता है;
और केवल $\True$ को निर्दिष्ट बनाने अथवा $\True$ तथा~$\Undef$ दोनों को
निर्दिष्ट बनाने का विकल्प भी है।

यहाँ हम सबसे प्रसिद्ध त्रिमूल्य तर्कशास्त्रों का चयन, उनकी प्रेरणाएँ और
उनके कुछ गुण प्रस्तुत करते हैं।

\end{document}
